\cotefile{83}
\scanpage{145}
\[
 k+k.v+k.w
\]
qui est de rang $3$, et pas égal à $P\ldots$

Je voudrais expliciter tout dans le cas où $H$ et $H'$ sont unipotentes, plus précisément
\[
 H=k\oplus L,\quad H'=k\oplus L'
\]
\margin{i.e. $\unsure{X|H},X'$ sont définis comme les schémas \deleted{\unsure{noy.}}\unsure{aux|vx} de $L\longrightarrow L^{\otimes2}$, $L'\longrightarrow L'^{\otimes2}$ élévation au carré}
\inserted{$P=H\otimes H'=\underbrace{k\oplus L\oplus L'}_{B}+L\otimes L'$}
avec $L,L'$ formés d'éléments de carré nul.

donc \deleted{$\unsure{D|T}_{H,H'}$} $T_{H,H'}$ s'identifie \deleted{\unsure{avec|à}} \inserted{à} $(L\otimes L')^\vee$. \deleted{Choisiss. \unsure{la|base} $u\in L,v\in L'$,} \inserted{La fonction}
\diagram{\begin{tikzcd}[ampersand replacement=\&,row sep=small]
\delta_{H,H'}:\ \& T_{H,H'} \arrow[r] \& \deleted{\unsure{k}}(L\otimes L')^{\otimes-2}\\
\& (L\otimes L')^{-1} \arrow[u,phantom,"\rotatebox{90}{$\simeq$}" description] \arrow[r,Rightarrow,swap,"\substack{\text{élévation}\\\text{\deleted{au} carré}}"] \& (L\otimes L')^{\otimes-2}
\end{tikzcd}}

donc
\[
 H\wedge H'\simeq k\oplus(L\otimes L')
\]
\margin{$L\otimes L'$ formé d'éléments de carré nul}

\deleted{Si $u\in L,v\in L'$}

L'hom
\[
 \deleted{\unsure{H}}X\times_S\deleted{\unsure{H'}}X'\longrightarrow X''=\deleted{\unsure{H}}X\wedge\deleted{\unsure{H}}X'
\]
est donné par
\[
 \underset{\substack{\text{sections}\\\text{de }L,L'\\\text{de carré nul}}}{(\xi,\eta)}
 \longrightarrow
 \underset{\substack{\text{section}\\\text{de }L\otimes L',\\\text{de carré}\\\text{nul}}}{\xi\otimes\eta}
\]
L'hom d'Alg correspondantes
\[
 H\otimes H'\simeq k+L+L'+L\otimes L'\longleftarrow H''=k\oplus L\otimes L'
\]
est celui défini par
% archive sheet number: 144

